\chapter{SPB Fitting}

\section{Objective}

The single parabolic band (SPB) fitting tools provide a bridge between measured
transport data and interpretable transport parameters. In TEDEX, the SPB
workflow is not one monolithic command. It is split into three complementary
entry points:

\begin{enumerate}[leftmargin=*]
  \item \code{effective_mass_fit.py}: fit density-of-states effective mass from
  Hall carrier concentration and Seebeck data.
  \item \code{performance_fit.py}: use Hall concentration, Seebeck, power
  factor, and optionally $zT$ to model power-factor and $zT$ trends.
  \item \code{conductivity_fit.py}: fit weighted mobility on a conductivity
  axis and plot Seebeck, power factor, and optionally $zT$ versus conductivity.
\end{enumerate}

The practical goal is to answer three questions: what effective mass is
consistent with the measured Pisarenko trend, which carrier concentration or
conductivity regime maximizes performance, and whether the measured data sit on
a physically plausible SPB curve.

\section{SPB Model Used in the Code}

The current implementation uses a compact acoustic-phonon SPB model. The
reduced Fermi level $\eta$ is obtained from the measured Seebeck magnitude. For
the default scattering exponent, the Seebeck magnitude is evaluated as

\begin{equation}
  |S| = \frac{k_B}{e}
  \left( \frac{2F_1(\eta)}{F_0(\eta)} - \eta \right),
\end{equation}

where $F_j(\eta)$ denotes the unnormalized Fermi integral

\begin{equation}
  F_j(\eta) = \int_0^\infty
  \frac{x^j}{1 + \exp(x-\eta)}\,dx.
\end{equation}

For a density-of-states effective mass $m^*$, the true carrier concentration is
computed as

\begin{equation}
  n =
  4\pi
  \left(
    \frac{2m^*m_e k_BT}{h^2}
  \right)^{3/2}
  F_{1/2}(\eta).
\end{equation}

By default, the Hall concentration is modeled as

\begin{equation}
  n_H = \frac{n}{r_H},
\end{equation}

where $r_H$ is the acoustic-SPB Hall factor. The command-line option
\code{--no-hall-factor} switches to the simplified convention $n_H=n$.

\section{Inputs and Outputs}

\begin{table}[htbp]
\centering
\caption{SPB fitting modes and their expected files.}
\begin{tabular}{@{}p{0.20\linewidth}p{0.32\linewidth}p{0.36\linewidth}@{}}
\toprule
Mode & Required data & Main outputs \\
\midrule
Effective mass & \code{nH}, \code{S}; optional \code{curve_id}, \code{T} &
\code{fit_summary.json}, \code{effective_mass_points.csv},
\code{pisarenko_curve.csv}, \code{pisarenko_fit.png/pdf} \\
Performance & \code{nH}, \code{S}, \code{PF}; optional \code{ZT}, parameters &
\code{performance_summary.json}, \code{performance_points.csv},
\code{performance_curve.csv}, \code{pf_fit.png/pdf}, \code{zt_fit.png/pdf} \\
Conductivity axis & \code{conductivity}, \code{S}; optional \code{PF},
\code{ZT}, \code{kL} & \code{conductivity_summary.json},
\code{conductivity_points.csv}, \code{conductivity_curve.csv},
\code{conductivity_*_fit.png/pdf} \\
\bottomrule
\end{tabular}
\end{table}

\section{Effective-Mass Fitting}

This mode is the first pass for a Hall-Pisarenko analysis. The script reads a
CSV containing Hall concentration and Seebeck columns, solves $\eta$ from
Seebeck, and fits the effective mass that best maps the SPB carrier
concentration curve to the measured $n_H$ values.


Single-sample fitting will automatically calculate average density-of-states effective mass based on provided data. We can manually 
define the effective mass by \code{--m 2.3} and alter fitting range by \code{--xlim 1e20 2e21}.

\twofigureslot
  {figures/chp4/4.1a.png}
  {(a) Automatic effective-mass fitting result.}
  {figures/chp4/4.1b.png}
  {(b) Manual effective-mass fitting result.}
  {Pisarenko fitting curves for Cu$_{2}$Se at 300 K.}

When multiple samples are required for fitting, we can use the \code{--group-by curve_id} option to group the samples 
by \code{curve_id}. This requires an original column named \code{curve_id} to be present in the input data. The paramter CSV is also need
to identify the fitting temperature if not at default 300 K.

\twofigureslot
  {figures/chp4/4.2a.png}
  {(a) Fitting of Ag2Se and Ag2SeTe at 300 K.}
  {figures/chp4/4.2b.pdf}
  {(b) Fitting of Cu2Se at 300 K and 750 K.}
  {Multi-groupPisarenko fitting curves.}

\begin{lstlisting}[language=bash,caption={Command slot: SPB effective-mass fit.}]
python main.py spb effective-mass \
  data/demo/spb_fitting/Ag2Se/Ag2Se_series.csv \
  --group-by curve_id \
  --params data/demo/spb_fitting/Ag2Se/Ag2Se_params.csv \
  --output-dir results/spb_fitting/effective_mass_fit \
  --formats png pdf \
  --legend outside \
  --no-show
\end{lstlisting}

\begin{lstlisting}[caption={Terminal-output slot: effective-mass fit.}]
# TODO: paste key saved-file lines here.
# Expected outputs include:
# results/spb_fitting/effective_mass_fit/fit_summary.json
# results/spb_fitting/effective_mass_fit/effective_mass_points.csv
# results/spb_fitting/effective_mass_fit/pisarenko_curve.csv
# results/spb_fitting/effective_mass_fit/pisarenko_fit.png
\end{lstlisting}


\section{Performance Fitting}

The performance fit extends the Pisarenko analysis by fitting or using a
mobility model and then rendering power-factor and $zT$ curves. This is useful
when the question is not only ``what $m^*$ explains $S v.s. (n_H)$?'' but also
``where should performance peak as carrier concentration changes?''

The performance fitting also supports single data, which just requires input fitting paramters in CLI line or by the computer.
If multiple curves is necessary, it is recommaned to use paramers.csv to input fittign parameters.

The Fig4.3 is an example of auto-fitting of weighted mobiliy to calculate $PF$ and $ZT$ under $\kappa_L = 0.44$ at fixed temeprature 750 K and
effective mass of 3.2 from Fig4.2b.

\twofigureslot
  {figures/chp4/4.3a.pdf}
  {(a) Fitting of $PF$ at 750 K.}
  {figures/chp4/4.3b.pdf}
  {(b) Fitting of $ZT$ at 750 K.}
  {Performance fitting of Cu2Se.}

To fit multipe sets of data, using Listing 4.3 with raw data and parameters, the results can be seen in Fig 4.4.

\twofigureslot{figures/chp4/4.4a.pdf}{(a) Fitting of $PF$ at 300 k.}{figures/chp4/4.4b.pdf}{(b) Fitting of $ZT$ at 300 K.}{Multiple sets data fitting of Ag2SSeTe.}

\begin{lstlisting}[language=bash,caption={Command slot: SPB performance fit.}]
python main.py spb performance \
  data/demo/spb_fitting/Ag2Se/Ag2Se_series.csv \
  --group-by curve_id \
  --params data/demo/spb_fitting/Ag2Se/Ag2Se_params.csv \
  --output-dir results/spb_fitting/performance_fit \
  --properties pf zt \
  --formats png pdf \
  --legend outside \
  --no-show
\end{lstlisting}

\begin{lstlisting}[caption={Terminal-output slot: performance fit.}]
# TODO: paste key saved-file lines here.
# Expected outputs include:
# results/spb_fitting/performance_fit/performance_summary.json
# results/spb_fitting/performance_fit/performance_curve.csv
# results/spb_fitting/performance_fit/pf_fit.png
# results/spb_fitting/performance_fit/zt_fit.png
\end{lstlisting}

\section{Conductivity-Axis Fitting}

The conductivity-axis workflow is useful when the experimental series is more
naturally interpreted as Seebeck, power factor, or $zT$ versus electrical
conductivity. The script fits or uses a weighted mobility prefactor and renders
model curves directly on the conductivity axis.

\begin{lstlisting}[language=bash,caption={Command slot: conductivity-axis SPB fit.}]
python main.py spb conductivity \
  data/demo/spb_fitting/Cu2SSeTe/900K.csv \
  --temperature 900 \
  --x conductivity \
  --y S \
  --pf PF \
  --zt ZT \
  --sigma-unit S/cm \
  --kL 0.3 \
  --output-dir results/spb_fitting/conductivity_fit \
  --properties seebeck pf zt \
  --formats png pdf \
  --legend outside \
  --no-show
\end{lstlisting}

\begin{lstlisting}[caption={Terminal-output slot: conductivity-axis fit.}]
# TODO: paste key saved-file lines here.
# Expected outputs include:
# results/spb_fitting/conductivity_fit/conductivity_summary.json
# results/spb_fitting/conductivity_fit/conductivity_curve.csv
# results/spb_fitting/conductivity_fit/conductivity_seebeck_fit.png
# results/spb_fitting/conductivity_fit/conductivity_pf_fit.png
# results/spb_fitting/conductivity_fit/conductivity_zt_fit.png
\end{lstlisting}

\twofigureslot{figures/chp4/4.5a.pdf}{(a) Fiteed $S$.}{figures/chp4/4.5b.pdf}{(b) Fitted $PF$.}{Fitted conducvtivity $v.s.$ thermoelectric parameetrs.}

\figureslot{figures/chp4/4.6.pdf}
{FFitted $ZT$ value.}

\section{Interpreting the Outputs}

\paragraph{Summary JSON files.}
The JSON summaries are the best place to check fitted scalar parameters such as
effective mass, mobility prefactor, lattice thermal conductivity setting, and
fit configuration. They should be archived with any figure used in a report.

\paragraph{Point CSV files.}
The point files keep the experimental measurements after unit normalization and
after derived SPB quantities have been added. These are useful for diagnosing
outliers, sign conventions, and incorrect input columns.

\paragraph{Curve CSV files.}
The curve files store the model line sampled over the plotting range. They make
the figure reproducible even if plotting styles change later.

\paragraph{Figures.}
The PNG/PDF outputs are intended for quick inspection and report insertion.
Prefer PDF or SVG-like vector outputs when the figure will be edited or placed
in a manuscript.

\section{Quality-Control Checklist}

\begin{enumerate}[leftmargin=*]
  \item Confirm the carrier concentration unit: \code{cm^-3} and \code{m^-3}
  differ by $10^6$.
  \item Confirm the Seebeck unit: \code{uV/K} and \code{V/K} differ by $10^6$.
  \item Confirm whether the Hall factor should be included. The default is to
  include the acoustic-SPB Hall factor.
  \item Avoid mixing temperatures in a single fit unless \code{--group-by} and
  \code{--T-column} or a parameter CSV explicitly separate them.
  \item Inspect residuals or point-vs-curve disagreement before interpreting a
  fitted effective mass as a material constant.
  \item Record the exact command, parameter CSV, and output directory for each
  figure used in a presentation or report.
\end{enumerate}

\section{Open Slots for Future Expansion}

\slotbox{Theory slot}{Add a fuller derivation of the SPB transport coefficients,
including Lorenz number and the selected scattering exponent.}

\slotbox{Validation slot}{Add one benchmark table comparing fitted values
against a literature SPB example or a manually reproduced calculation.}

\slotbox{Lab-data slot}{Add a CHY batch example once a complete Hall-containing
dataset is selected for publication-style analysis.}
